% ============================================================
%  C: Como Programar -- 6ª Edição | Deitel & Deitel
%  Capítulo 4 -- Controle de Programa em C
%  EXERCÍCIOS (4.5 -- 4.37)
%  Abordagem incremental: Caps. 1 + 2 + 3 + 4
% ============================================================
\documentclass[12pt,a4paper]{article}
\usepackage[utf8]{inputenc}
\usepackage[T1]{fontenc}
\usepackage[brazil]{babel}
\usepackage{geometry}
\geometry{top=2.5cm,bottom=2.5cm,left=3cm,right=3cm}
\usepackage{parskip}\setlength{\parskip}{6pt}\setlength{\parindent}{0pt}
\usepackage{xcolor,tcolorbox,enumitem,listings,fancyhdr,titlesec,hyperref,booktabs,array,amsmath}
\tcbuselibrary{skins,breakable}

\definecolor{deitelblue}{RGB}{0,70,127}
\definecolor{deitelgray}{RGB}{240,242,245}
\definecolor{deitelgreen}{RGB}{0,110,60}
\definecolor{deitellightblue}{RGB}{220,235,250}
\definecolor{deitelred}{RGB}{160,20,20}
\definecolor{deitellorange}{RGB}{200,90,0}

\newtcolorbox{boaprática}[1][]{colback=deitellightblue,colframe=deitelblue,
  fonttitle=\bfseries\small,title={Boa prática de programação},breakable,#1}
\newtcolorbox{erroco}[1][]{colback=red!5,colframe=deitelred,
  fonttitle=\bfseries\small,title={Erro comum de programação},breakable,#1}
\newtcolorbox{respostabox}[1][]{colback=deitelgray,colframe=deitelblue!60,
  fonttitle=\bfseries\small\color{deitelblue},title={Resposta detalhada},breakable,#1}
\newtcolorbox{enunciadoBox}[1][]{colback=white,colframe=deitelblue,
  fonttitle=\bfseries,title={#1},breakable}
\newtcolorbox{saida}[1][]{colback=black!88,colframe=deitelblue,
  fontupper=\ttfamily\small\color{green!70},
  title={\small\color{white}Saída do programa},breakable,#1}

\setlist[enumerate,1]{label=\textbf{\alph*)},leftmargin=2em}
\setlist[itemize]{leftmargin=2em}

\lstset{
  basicstyle=\ttfamily\small,backgroundcolor=\color{deitelgray},
  frame=single,framerule=0.4pt,rulecolor=\color{deitelblue!40},
  breaklines=true,showstringspaces=false,
  keywordstyle=\color{deitelblue}\bfseries,
  commentstyle=\color{deitelgreen}\itshape,
  stringstyle=\color{deitelred},
  language=C,numbers=left,numberstyle=\tiny\color{gray},
  xleftmargin=18pt,xrightmargin=5pt,stepnumber=1,
}

\pagestyle{fancy}\fancyhf{}
\fancyhead[L]{\small\color{deitelblue}\textbf{C: Como Programar -- 6ª ed. | Deitel \& Deitel}}
\fancyhead[R]{\small\color{deitelblue}Cap. 4 -- Exercícios}
\fancyfoot[C]{\small\thepage}
\renewcommand{\headrulewidth}{0.4pt}

\titleformat{\section}[block]{\large\bfseries\color{deitelblue}}{\thesection.}{0.5em}{}[\titlerule]
\titleformat{\subsection}[block]{\normalsize\bfseries\color{deitelblue!80}}{}{}{}

\hypersetup{colorlinks=true,linkcolor=deitelblue}

\begin{document}

\begin{titlepage}
  \centering\vspace*{2cm}
  {\color{deitelblue}\rule{\linewidth}{2pt}}\\[0.4cm]
  {\huge\bfseries\color{deitelblue}C: Como Programar}\\[0.2cm]
  {\large\color{deitelblue}6ª Edição $\cdot$ Paul Deitel \& Harvey Deitel}\\[0.2cm]
  {\color{deitelblue}\rule{\linewidth}{2pt}}\\[1cm]
  {\LARGE\bfseries Capítulo 4}\\[0.3cm]
  {\Large\color{deitelblue!80}Controle de Programa em C}\\[0.8cm]
  \begin{tcolorbox}[colback=deitellightblue,colframe=deitelblue,width=0.85\linewidth]
    \centering{\large\bfseries Exercícios (4.5--4.37)}\\[0.2cm]
    {\normalsize Resolvidos passo a passo, com código completo}
  \end{tcolorbox}
\end{titlepage}

\tableofcontents\newpage

% ============================================================
\section{Exercício 4.5 -- Identificar e Corrigir Erros}
% ============================================================

\begin{respostabox}
\begin{enumerate}[label=\textbf{\alph*)}]

\item \textbf{Código:} \texttt{For ( x = 100, x >= 1, x++ )}\\
\textbf{Erros:} (1) \texttt{For} com \texttt{F} maiúsculo (C é case-sensitive);
(2) usa-se \textit{vírgula} em vez de \textit{ponto e vírgula} entre as expressões;
(3) com \texttt{x++} e condição \texttt{x >= 1}, criamos um loop infinito (deveria
ser \texttt{x--}).\\
\textbf{Correção:}
\begin{lstlisting}
for ( x = 100; x >= 1; x-- )
    printf( "%d\n", x );
\end{lstlisting}

\item \textbf{Switch para par/ímpar:} Falta \texttt{break} no \texttt{case 0}.\\
\textbf{Correção:}
\begin{lstlisting}
switch ( valor % 2 ) {
    case 0:
        printf( "Inteiro par\n" );
        break;     /* adiciona break! */
    case 1:
        printf( "Inteiro impar\n" );
        break;
}
\end{lstlisting}

\item \textbf{Leitura de inteiro e caractere:} O \texttt{scanf("\%d", \&intVal)} deixa
o caractere \texttt{'\textbackslash n'} (Enter) no buffer, e \texttt{getchar} captura
esse \texttt{'\textbackslash n'} em vez do \texttt{'A'}.\\
\textbf{Correção:} consumir o caractere extra antes do \texttt{getchar}.
\begin{lstlisting}
scanf( "%d", &intVal );
scanf( " %c", &charVal );  /* o espaco antes de %c pula brancos */
printf( "Inteiro: %d\nCaractere: %c\n", intVal, charVal );
\end{lstlisting}

\item \textbf{Ponto flutuante em \texttt{for}:} O mesmo problema de imprecisão do
Exercício 4.4(b). A condição \texttt{x == .0001} provavelmente nunca será
verdadeira exatamente.\\
\textbf{Correção:} usar inteiro como variável de controle.

\item \textbf{Inteiros ímpares de 999 a 1:} \texttt{x += 2} \textit{aumenta} \texttt{x},
mas a intenção é \textit{diminuir} (de 999 para 1).\\
\textbf{Correção:}
\begin{lstlisting}
for ( x = 999; x >= 1; x -= 2 ) {
    printf( "%d\n", x );
}
\end{lstlisting}

\item \textbf{\texttt{do...while} com erros de capitalização:} \texttt{Do} e \texttt{While}
devem ser minúsculos. Além disso, a condição \texttt{< 100} exclui o 100 (o último
par desejado); deve ser \texttt{<=}.\\
\textbf{Correção:}
\begin{lstlisting}
contador = 2;
do {
    if ( contador % 2 == 0 ) {
        printf( "%d\n", contador );
    }
    contador += 2;
} while ( contador <= 100 );
\end{lstlisting}

\item \textbf{Soma de 100 a 150:} O \texttt{;} após \texttt{)} cria um corpo vazio
para o \texttt{for}; \texttt{total += x} executa apenas uma vez ao final.\\
\textbf{Correção:} remover o \texttt{;}.
\begin{lstlisting}
for ( x = 100; x <= 150; x++ ) {  /* sem ; depois do ) */
    total += x;
}
\end{lstlisting}

\end{enumerate}
\end{respostabox}

\newpage

% ============================================================
\section{Exercício 4.6 -- Valores impressos por estruturas \texttt{for}}
% ============================================================

\begin{respostabox}
\begin{enumerate}[label=\textbf{\alph*)}]

\item \texttt{for ( x = 2; x <= 13; x += 2 )} \quad
$\to$ \textbf{2, 4, 6, 8, 10, 12} (não 14, pois $14 > 13$)

\item \texttt{for ( x = 5; x <= 22; x += 7 )} \quad
$\to$ \textbf{5, 12, 19} (próximo seria 26, mas $26 > 22$)

\item \texttt{for ( x = 3; x <= 15; x += 3 )} \quad
$\to$ \textbf{3, 6, 9, 12, 15}

\item \texttt{for ( x = 1; x <= 5; x += 7 )} \quad
$\to$ \textbf{1} apenas (próximo seria 8, mas $8 > 5$)

\item \texttt{for ( x = 12; x >= 2; x -= 3 )} \quad
$\to$ \textbf{12, 9, 6, 3} (próximo seria 0, mas $0 < 2$)

\end{enumerate}
\end{respostabox}

% ============================================================
\section{Exercício 4.7 -- Estruturas \texttt{for} para sequências}
% ============================================================

\begin{respostabox}
\begin{enumerate}[label=\textbf{\alph*)}]

\item \textbf{1, 2, 3, 4, 5, 6, 7:}
\begin{lstlisting}
for ( i = 1; i <= 7; ++i ) {
    printf( "%d ", i );
}
\end{lstlisting}

\item \textbf{3, 8, 13, 18, 23} (incremento de 5):
\begin{lstlisting}
for ( i = 3; i <= 23; i += 5 ) {
    printf( "%d ", i );
}
\end{lstlisting}

\item \textbf{20, 14, 8, 2, $-$4, $-$10} (decremento de 6):
\begin{lstlisting}
for ( i = 20; i >= -10; i -= 6 ) {
    printf( "%d ", i );
}
\end{lstlisting}

\item \textbf{19, 27, 35, 43, 51} (incremento de 8):
\begin{lstlisting}
for ( i = 19; i <= 51; i += 8 ) {
    printf( "%d ", i );
}
\end{lstlisting}

\end{enumerate}
\end{respostabox}

% ============================================================
\section{Exercício 4.8 -- O que o programa faz?}
% ============================================================

\begin{respostabox}
O programa lê dois inteiros \texttt{x} e \texttt{y} (entre 1 e 20). Em seguida, usa
\textbf{loops aninhados}: o externo executa \texttt{y} vezes (cada iteração imprime
uma linha) e o interno executa \texttt{x} vezes (imprimindo \texttt{@} em cada).

\textbf{Resultado:} um \textbf{retângulo de arrobas} com \texttt{y} linhas e
\texttt{x} colunas.

\textbf{Exemplo:} se \texttt{x = 5, y = 3}, a saída será:

\begin{saida}
@@@@@\\
@@@@@\\
@@@@@
\end{saida}
\end{respostabox}

\newpage

% ============================================================
\section{Exercício 4.9 -- Soma de uma Sequência de Inteiros}
% ============================================================

\begin{respostabox}
\begin{lstlisting}
/* Ex4.9: soma_sequencia.c
   Le N (primeiro valor), depois N inteiros, e exibe a soma */
#include <stdio.h>

int main( void )
{
    int n;        /* numero de valores a somar */
    int valor;    /* valor lido a cada iteracao */
    int soma = 0; /* total acumulado            */
    int i;

    printf( "Digite quantos valores voce quer somar: " );
    scanf( "%d", &n );

    for ( i = 1; i <= n; ++i ) {
        printf( "Digite o valor %d: ", i );
        scanf( "%d", &valor );
        soma += valor;
    }

    printf( "A soma e: %d\n", soma );
    return 0;
} /* fim da funcao main */
\end{lstlisting}
\end{respostabox}

% ============================================================
\section{Exercício 4.10 -- Média de uma Sequência (Sentinela 9999)}
% ============================================================

\begin{respostabox}
\begin{lstlisting}
/* Ex4.10: media_sentinela.c
   Calcula a media usando 9999 como sentinela */
#include <stdio.h>

int main( void )
{
    int valor;
    int soma = 0;
    int contador = 0;
    float media;

    printf( "Digite inteiros (9999 para terminar): " );
    scanf( "%d", &valor );

    while ( valor != 9999 ) {
        soma += valor;
        ++contador;
        scanf( "%d", &valor );
    }

    if ( contador != 0 ) {
        media = (float) soma / contador;  /* conversao para float */
        printf( "Media de %d valores: %.2f\n", contador, media );
    }
    else {
        printf( "Nenhum valor foi informado.\n" );
    }

    return 0;
} /* fim da funcao main */
\end{lstlisting}

\textbf{Conversão de tipo (cast):} \texttt{(float) soma} converte \texttt{soma} para
\texttt{float} antes da divisão, garantindo divisão de ponto flutuante (com casas
decimais), em vez de divisão inteira (que truncaria).
\end{respostabox}

% ============================================================
\section{Exercício 4.11 -- Achar o Menor}
% ============================================================

\begin{respostabox}
\begin{lstlisting}
/* Ex4.11: menor.c
   Le N, depois N inteiros, e exibe o menor */
#include <stdio.h>

int main( void )
{
    int n, valor, menor, i;

    printf( "Digite quantos valores comparar: " );
    scanf( "%d", &n );

    printf( "Digite o 1o valor: " );
    scanf( "%d", &menor );  /* primeiro valor inicializa menor */

    for ( i = 2; i <= n; ++i ) {
        printf( "Digite o %do valor: ", i );
        scanf( "%d", &valor );
        if ( valor < menor ) {
            menor = valor;
        }
    }

    printf( "O menor valor e: %d\n", menor );
    return 0;
} /* fim da funcao main */
\end{lstlisting}
\end{respostabox}

% ============================================================
\section{Exercício 4.12 -- Soma dos Pares de 2 a 30}
% ============================================================

\begin{respostabox}
\begin{lstlisting}
/* Ex4.12: soma_pares.c
   Soma os inteiros pares de 2 a 30 */
#include <stdio.h>

int main( void )
{
    int soma = 0;
    int i;

    for ( i = 2; i <= 30; i += 2 ) {
        soma += i;
    }

    printf( "Soma dos pares de 2 a 30: %d\n", soma );
    return 0;
} /* fim da funcao main */
\end{lstlisting}

\textbf{Resultado:} $2+4+6+\ldots+30 = 240$.
\end{respostabox}

% ============================================================
\section{Exercício 4.13 -- Produto dos Ímpares de 1 a 15}
% ============================================================

\begin{respostabox}
\begin{lstlisting}
/* Ex4.13: produto_impares.c
   Produto dos inteiros impares de 1 a 15 */
#include <stdio.h>

int main( void )
{
    int produto = 1;  /* atencao: produto inicia em 1, nao 0 */
    int i;

    for ( i = 1; i <= 15; i += 2 ) {
        produto *= i;
    }

    printf( "Produto dos impares de 1 a 15: %d\n", produto );
    return 0;
} /* fim da funcao main */
\end{lstlisting}

\textbf{Resultado:} $1 \times 3 \times 5 \times 7 \times 9 \times 11 \times 13 \times 15
= 2.027.025$.

\begin{boaprática}
  \textbf{Atenção à inicialização:} totais começam com \texttt{0}, mas produtos
  com \texttt{1} (elemento neutro da multiplicação). Inicializar produto com 0
  é um erro comum -- todo o resultado seria 0!
\end{boaprática}
\end{respostabox}

\newpage

% ============================================================
\section{Exercício 4.14 -- Tabela de Fatoriais (1 a 5)}
% ============================================================

\begin{respostabox}
\begin{lstlisting}
/* Ex4.14: fatoriais.c
   Imprime tabela de fatoriais de 1 a 5 */
#include <stdio.h>

int main( void )
{
    int n, i;
    long fatorial;

    printf( "%5s%15s\n", "n", "n!" );

    for ( n = 1; n <= 5; ++n ) {
        fatorial = 1;
        for ( i = 1; i <= n; ++i ) {
            fatorial *= i;
        }
        printf( "%5d%15ld\n", n, fatorial );
    }

    return 0;
} /* fim da funcao main */
\end{lstlisting}

\begin{saida}
\hphantom{xx}n\hphantom{xxxxxxxxxxxx}n!\\
\hphantom{xx}1\hphantom{xxxxxxxxxxxxx}1\\
\hphantom{xx}2\hphantom{xxxxxxxxxxxxx}2\\
\hphantom{xx}3\hphantom{xxxxxxxxxxxxx}6\\
\hphantom{xx}4\hphantom{xxxxxxxxxxxx}24\\
\hphantom{xx}5\hphantom{xxxxxxxxxxx}120
\end{saida}

\textbf{Por que dificuldade em $20!$?} O fatorial cresce \textit{explosivamente}:
\[
  20! = 2\,432\,902\,008\,176\,640\,000 \approx 2{,}4 \times 10^{18}
\]
Esse valor ultrapassa o limite de \texttt{long} em sistemas de 32 bits ($\approx 2{,}1 \times 10^9$).
Mesmo em sistemas de 64 bits, $21! \approx 5{,}1 \times 10^{19}$ já excede o limite de
\texttt{long long} (\(\approx 9{,}2 \times 10^{18}\)). Para fatoriais grandes, é preciso
usar \texttt{double} (com perda de precisão para valores exatos) ou bibliotecas de
aritmética arbitrária.
\end{respostabox}

% ============================================================
\section{Exercício 4.15 -- Juros Compostos com Várias Taxas}
% ============================================================

\begin{respostabox}
\begin{lstlisting}
/* Ex4.15: juros_compostos_taxas.c
   Tabela de juros compostos para taxas de 5% a 10% */
#include <stdio.h>
#include <math.h>

int main( void )
{
    double principal = 1000.00;  /* valor inicial */
    double valor;
    int taxa, ano;

    for ( taxa = 5; taxa <= 10; ++taxa ) {
        printf( "\n=== Taxa de %d%% ===\n", taxa );
        printf( "%5s%21s\n", "Ano", "Valor acumulado" );

        for ( ano = 1; ano <= 10; ++ano ) {
            valor = principal * pow( 1.0 + taxa / 100.0, ano );
            printf( "%5d%21.2f\n", ano, valor );
        }
    }

    return 0;
} /* fim da funcao main */
\end{lstlisting}

\textbf{Fórmula de juros compostos:} $V = P \cdot (1 + r)^n$, onde $P$ é o principal,
$r$ a taxa em decimal e $n$ o número de anos.
\end{respostabox}

\newpage

% ============================================================
\section{Exercício 4.16 -- Quatro Padrões de Triângulos}
% ============================================================

\begin{respostabox}
\begin{lstlisting}
/* Ex4.16: triangulos.c
   Imprime 4 padroes triangulares com loops for aninhados */
#include <stdio.h>

int main( void )
{
    int linha, j;

    /* (A) Triangulo crescente alinhado a esquerda */
    printf( "(A)\n" );
    for ( linha = 1; linha <= 10; ++linha ) {
        for ( j = 1; j <= linha; ++j ) {
            printf( "*" );
        }
        printf( "\n" );
    }

    /* (B) Triangulo decrescente alinhado a esquerda */
    printf( "\n(B)\n" );
    for ( linha = 10; linha >= 1; --linha ) {
        for ( j = 1; j <= linha; ++j ) {
            printf( "*" );
        }
        printf( "\n" );
    }

    /* (C) Triangulo decrescente alinhado a direita */
    printf( "\n(C)\n" );
    for ( linha = 10; linha >= 1; --linha ) {
        for ( j = 1; j <= 10 - linha; ++j ) {
            printf( " " );    /* espacos antes  */
        }
        for ( j = 1; j <= linha; ++j ) {
            printf( "*" );
        }
        printf( "\n" );
    }

    /* (D) Triangulo crescente alinhado a direita */
    printf( "\n(D)\n" );
    for ( linha = 1; linha <= 10; ++linha ) {
        for ( j = 1; j <= 10 - linha; ++j ) {
            printf( " " );
        }
        for ( j = 1; j <= linha; ++j ) {
            printf( "*" );
        }
        printf( "\n" );
    }

    return 0;
} /* fim da funcao main */
\end{lstlisting}
\end{respostabox}

% ============================================================
\section{Exercício 4.17 -- Limites de Crédito Cortados pela Metade}
% ============================================================

\begin{respostabox}
\begin{lstlisting}
/* Ex4.17: limites_credito_metade.c
   Analisa 3 clientes apos corte de 50% no limite de credito */
#include <stdio.h>

int main( void )
{
    int conta;
    float limiteAntigo, limiteNovo, saldo;
    int cliente;

    for ( cliente = 1; cliente <= 3; ++cliente ) {
        printf( "\n=== Cliente %d ===\n", cliente );
        printf( "Numero da conta: " );
        scanf( "%d", &conta );

        printf( "Limite de credito antigo: " );
        scanf( "%f", &limiteAntigo );

        printf( "Saldo atual: " );
        scanf( "%f", &saldo );

        limiteNovo = limiteAntigo / 2.0f;

        printf( "Conta: %d\n", conta );
        printf( "Limite antigo: R$ %.2f\n", limiteAntigo );
        printf( "Limite novo:   R$ %.2f\n", limiteNovo );
        printf( "Saldo atual:   R$ %.2f\n", saldo );

        if ( saldo > limiteNovo ) {
            printf( ">>> Saldo excede o novo limite! <<<\n" );
        }
        else {
            printf( "Saldo dentro do novo limite.\n" );
        }
    }

    return 0;
} /* fim da funcao main */
\end{lstlisting}
\end{respostabox}

\newpage

% ============================================================
\section{Exercício 4.18 -- Gráfico de Barras (Histograma)}
% ============================================================

\begin{respostabox}
\begin{lstlisting}
/* Ex4.18: grafico_barras.c
   Le 5 numeros (1 a 30) e exibe linha de asteriscos correspondente */
#include <stdio.h>

int main( void )
{
    int i, j, valor;

    printf( "Digite 5 numeros entre 1 e 30:\n" );

    for ( i = 1; i <= 5; ++i ) {
        scanf( "%d", &valor );

        for ( j = 1; j <= valor; ++j ) {
            printf( "*" );
        }
        printf( "\n" );
    }

    return 0;
} /* fim da funcao main */
\end{lstlisting}

\textbf{Exemplo de execução:}
\begin{saida}
Digite 5 numeros entre 1 e 30:\\
3 7 12 5 9\\
***\\
*******\\
************\\
*****\\
*********
\end{saida}
\end{respostabox}

% ============================================================
\section{Exercício 4.19 -- Calculando Vendas com \texttt{switch}}
% ============================================================

\begin{respostabox}
\begin{lstlisting}
/* Ex4.19: vendas_switch.c
   Calcula valor total de vendas usando switch */
#include <stdio.h>

int main( void )
{
    int produto, quantidade;
    float preco, total = 0.0f;

    printf( "Digite numero do produto e quantidade vendida\n" );
    printf( "(produto = -1 para terminar):\n" );

    scanf( "%d", &produto );

    while ( produto != -1 ) {
        scanf( "%d", &quantidade );

        switch ( produto ) {
            case 1:
                preco = 2.98f;
                break;
            case 2:
                preco = 4.50f;
                break;
            case 3:
                preco = 9.98f;
                break;
            case 4:
                preco = 4.49f;
                break;
            case 5:
                preco = 6.87f;
                break;
            default:
                preco = 0.0f;
                printf( "Produto invalido: %d\n", produto );
                break;
        }

        total += preco * quantidade;

        printf( "Proximo produto (-1 para terminar): " );
        scanf( "%d", &produto );
    }

    printf( "\nValor total de vendas: R$ %.2f\n", total );
    return 0;
} /* fim da funcao main */
\end{lstlisting}
\end{respostabox}

\newpage

% ============================================================
\section{Exercício 4.20 -- Tabelas-Verdade}
% ============================================================

\begin{respostabox}

\textbf{Tabela do operador \texttt{\&\&} (E lógico):}
\begin{center}
\renewcommand{\arraystretch}{1.4}
\begin{tabular}{c c c}
  \toprule
  \textbf{Condição1} & \textbf{Condição2} & \textbf{Cond1 \&\& Cond2} \\ \midrule
  0 & 0 & 0 \\
  0 & não zero & 0 \\
  não zero & 0 & \textbf{0} \\
  não zero & não zero & \textbf{1} \\
  \bottomrule
\end{tabular}
\end{center}

\medskip\textbf{Tabela do operador \texttt{||} (OU lógico):}
\begin{center}
\renewcommand{\arraystretch}{1.4}
\begin{tabular}{c c c}
  \toprule
  \textbf{Condição1} & \textbf{Condição2} & \textbf{Cond1 || Cond2} \\ \midrule
  0 & 0 & 0 \\
  0 & não zero & 1 \\
  não zero & 0 & \textbf{1} \\
  não zero & não zero & \textbf{1} \\
  \bottomrule
\end{tabular}
\end{center}

\medskip\textbf{Tabela do operador \texttt{!} (NÃO lógico):}
\begin{center}
\renewcommand{\arraystretch}{1.4}
\begin{tabular}{c c}
  \toprule
  \textbf{Condição1} & \textbf{!Condição1} \\ \midrule
  0 & 1 \\
  não zero & \textbf{0} \\
  \bottomrule
\end{tabular}
\end{center}

\textbf{Em C:} qualquer valor diferente de zero é tratado como ``verdadeiro''
(representado por 1 na saída de operadores lógicos), e zero é ``falso''.
\end{respostabox}

% ============================================================
\section{Exercício 4.21 -- Inicialização da Variável de Controle na Declaração}
% ============================================================

\begin{respostabox}
A Figura~4.2 do livro usa \texttt{contador} como variável de controle e o inicializa
\textit{dentro} do cabeçalho do \texttt{for}. Aqui, fazemos a inicialização junto
com a declaração:

\begin{lstlisting}
/* Ex4.21: contador na declaracao */
#include <stdio.h>

int main( void )
{
    int contador = 1;  /* inicializa na declaracao */

    for ( ; contador <= 10; ++contador ) {  /* primeiro campo vazio */
        printf( "%d\n", contador );
    }

    return 0;
} /* fim da funcao main */
\end{lstlisting}

\textit{Note:} no cabeçalho \texttt{for}, o primeiro campo (inicialização) está vazio.
A inicialização foi feita na declaração da variável.
\end{respostabox}

\newpage

% ============================================================
\section{Exercício 4.22 -- Nota Média da Classe}
% ============================================================

\begin{respostabox}
\begin{lstlisting}
/* Ex4.22: nota_media_classe.c
   Modificacao do programa de notas-letra para calcular media */
#include <stdio.h>

int main( void )
{
    int nota;
    int total = 0;
    int contador = 0;
    float media;

    printf( "Digite as notas (EOF para terminar):\n" );

    while ( ( nota = getchar() ) != EOF ) {
        if ( nota >= '0' && nota <= '9' ) {
            /* Aqui simplesmente lemos numero por numero
               via scanf seria mais robusto */
        }
    }

    /* Versao simplificada com scanf: */
    printf( "Digite uma nota (-1 para terminar): " );
    scanf( "%d", &nota );

    while ( nota != -1 ) {
        total += nota;
        ++contador;
        printf( "Proxima nota (-1 para terminar): " );
        scanf( "%d", &nota );
    }

    if ( contador != 0 ) {
        media = (float) total / contador;
        printf( "Media da classe: %.2f\n", media );
    }
    else {
        printf( "Nenhuma nota informada.\n" );
    }

    return 0;
} /* fim da funcao main */
\end{lstlisting}
\end{respostabox}

% ============================================================
\section{Exercício 4.24 -- Avaliação de Expressões}
% ============================================================

\begin{enunciadoBox}[Enunciado 4.24]
Considere $i = 1$, $j = 2$, $k = 3$, $m = 2$. Determine o valor impresso por cada
\texttt{printf}.
\end{enunciadoBox}

\begin{respostabox}
Em C, expressões relacionais e lógicas avaliam para \textbf{1} (verdadeiro) ou
\textbf{0} (falso).

\begin{enumerate}[label=\textbf{\alph*)}]
\item \texttt{i == 1} $\to$ $1 == 1$ é \textbf{verdadeiro} $\to$ imprime \texttt{1}
\item \texttt{j == 3} $\to$ $2 == 3$ é \textbf{falso} $\to$ imprime \texttt{0}
\item \texttt{i >= 1 \&\& j < 4} $\to$ $1 \geq 1$ (V) E $2 < 4$ (V) $\to$ \texttt{1}
\item \texttt{m <= 99 \&\& k < m} $\to$ $2 \leq 99$ (V) E $3 < 2$ (F) $\to$ \texttt{0}
\item \texttt{j >= i || k == m} $\to$ $2 \geq 1$ (V) OU $3 == 2$ (F) $\to$ \texttt{1}
\item \texttt{k+m < j || 3-j >= k} $\to$ $5 < 2$ (F) OU $1 \geq 3$ (F) $\to$ \texttt{0}
\item \texttt{!m} $\to$ \texttt{!2} $\to$ \texttt{0} (qualquer não-zero \(\to\) 0)
\item \texttt{!(j - m)} $\to$ \texttt{!(2-2)} $\to$ \texttt{!0} $\to$ \texttt{1}
\item \texttt{!(k > m)} $\to$ \texttt{!(3 > 2)} $\to$ \texttt{!1} $\to$ \texttt{0}
\item \texttt{!(j > k)} $\to$ \texttt{!(2 > 3)} $\to$ \texttt{!0} $\to$ \texttt{1}
\end{enumerate}
\end{respostabox}

\newpage

% ============================================================
\section{Exercício 4.25 -- Tabela Decimal/Binário/Octal/Hex}
% ============================================================

\begin{respostabox}
\begin{lstlisting}
/* Ex4.25: tabela_bases.c
   Imprime equivalentes binario, octal e hex de 1 a 256 */
#include <stdio.h>

int main( void )
{
    int decimal, n, binario, valorPos;

    printf( "%10s%10s%10s%10s\n", "Decimal", "Binario", "Octal", "Hex" );

    for ( decimal = 1; decimal <= 256; ++decimal ) {
        /* Calcula representacao binaria como inteiro decimal */
        binario = 0;
        valorPos = 1;
        n = decimal;
        while ( n > 0 ) {
            binario += ( n % 2 ) * valorPos;
            n /= 2;
            valorPos *= 10;
        }

        /* Octal e hex impressos com %o e %X (suportados por printf) */
        printf( "%10d%10d%10o%10X\n", decimal, binario, decimal, decimal );
    }

    return 0;
} /* fim da funcao main */
\end{lstlisting}

\textbf{Observação:} \texttt{printf} oferece especificadores diretos para octal
(\texttt{\%o}) e hexadecimal (\texttt{\%X}). Para o binário, calculamos manualmente,
pois C não tem especificador embutido.
\end{respostabox}

% ============================================================
\section{Exercício 4.26 -- Cálculo de $\pi$ por Série Infinita}
% ============================================================

\begin{respostabox}
\textbf{Fórmula (Leibniz):}
$\pi \approx 4 - \dfrac{4}{3} + \dfrac{4}{5} - \dfrac{4}{7} + \dfrac{4}{9} - \cdots$

\begin{lstlisting}
/* Ex4.26: pi_serie.c
   Aproxima pi pela serie de Leibniz e exibe convergencia */
#include <stdio.h>
#include <math.h>

int main( void )
{
    double pi = 0.0;
    double termo;
    int sinal = 1;
    int n;
    int denominador;

    printf( "%10s%20s\n", "Termos", "Pi aproximado" );

    for ( n = 1; n <= 100000; ++n ) {
        denominador = 2 * n - 1;       /* 1, 3, 5, 7, 9, ... */
        termo = 4.0 / denominador;
        pi += sinal * termo;
        sinal = -sinal;                /* alterna o sinal */

        if ( n <= 10 || n == 100 || n == 1000 ||
             n == 10000 || n == 100000 ) {
            printf( "%10d%20.6f\n", n, pi );
        }
    }

    return 0;
} /* fim da funcao main */
\end{lstlisting}

\textbf{Análise da convergência (lenta!):}
\begin{itemize}[noitemsep]
  \item Após 10 termos: $\pi \approx 3{,}04$
  \item Após 100 termos: $\pi \approx 3{,}13$ (1 casa correta)
  \item Após 1.000 termos: $\pi \approx 3{,}1406$ (3 casas)
  \item Após 10.000 termos: $\pi \approx 3{,}14149$ (4 casas)
  \item Após 100.000 termos: $\pi \approx 3{,}14158$ (5 casas)
\end{itemize}

\textbf{Quantos termos para cada precisão?}
\begin{itemize}[noitemsep]
  \item Para 3,14: cerca de 200 termos
  \item Para 3,141: cerca de 1.000 termos
  \item Para 3,1415: cerca de 10.000 termos
  \item Para 3,14159: cerca de 100.000 termos
\end{itemize}
A série de Leibniz é \textbf{notoriamente lenta} -- cada casa decimal extra exige
aproximadamente 10× mais termos!
\end{respostabox}

\newpage

% ============================================================
\section{Exercício 4.27 -- Triplas de Pitágoras (Força Bruta)}
% ============================================================

\begin{respostabox}
\textbf{Critério:} $\text{cateto1}^2 + \text{cateto2}^2 = \text{hipotenusa}^2$.

\begin{lstlisting}
/* Ex4.27: triplas_pitagoras.c
   Encontra todas as triplas de Pitagoras com lados <= 500 */
#include <stdio.h>

int main( void )
{
    int cat1, cat2, hip;
    int contador = 0;

    printf( "%6s%6s%6s\n", "Cat1", "Cat2", "Hip" );

    for ( cat1 = 1; cat1 <= 500; ++cat1 ) {
        for ( cat2 = cat1; cat2 <= 500; ++cat2 ) {  /* cat2 >= cat1 */
            for ( hip = cat2; hip <= 500; ++hip ) {  /* hip >= cat2 */
                if ( cat1*cat1 + cat2*cat2 == hip*hip ) {
                    printf( "%6d%6d%6d\n", cat1, cat2, hip );
                    ++contador;
                }
            }
        }
    }

    printf( "\nTotal de triplas encontradas: %d\n", contador );
    return 0;
} /* fim da funcao main */
\end{lstlisting}

\textbf{Otimização:} para evitar duplicatas (3,4,5 = 4,3,5), exigimos
$\text{cat1} \leq \text{cat2} \leq \text{hip}$. As primeiras triplas:
(3,4,5), (5,12,13), (6,8,10), (7,24,25), (8,15,17), (9,12,15), (9,40,41), \ldots

\textbf{Reflexão sobre força bruta:} Este programa testa $500^3 = 125$ milhões de
combinações. CPUs modernas executam isso em frações de segundo. Esse exemplo
ilustra como o ``poder bruto'' computacional torna soluções triviais de problemas
que antes pareciam inviáveis.
\end{respostabox}

% ============================================================
\section{Exercício 4.28 -- Pagamento Semanal com Múltiplos Códigos}
% ============================================================

\begin{respostabox}
\begin{lstlisting}
/* Ex4.28: pagamento_semanal.c
   Calcula pagamento semanal para 4 tipos de funcionario */
#include <stdio.h>

int main( void )
{
    int codigo;
    float salario, pagamento, vendas, valorPeca;
    int horas, pecas;

    printf( "Codigos: 1=gerente, 2=horista, 3=comissao, 4=peca\n" );
    printf( "Codigo do funcionario (-1 para terminar): " );
    scanf( "%d", &codigo );

    while ( codigo != -1 ) {
        switch ( codigo ) {
            case 1:  /* Gerente */
                printf( "Salario semanal fixo: " );
                scanf( "%f", &salario );
                pagamento = salario;
                break;

            case 2:  /* Horista */
                printf( "Horas trabalhadas: " );
                scanf( "%d", &horas );
                printf( "Salario por hora: " );
                scanf( "%f", &salario );
                if ( horas <= 40 ) {
                    pagamento = horas * salario;
                }
                else {
                    pagamento = 40 * salario +
                                ( horas - 40 ) * salario * 1.5f;
                }
                break;

            case 3:  /* Comissao */
                printf( "Vendas brutas semanais: " );
                scanf( "%f", &vendas );
                pagamento = 250.0f + 0.057f * vendas;
                break;

            case 4:  /* Por peca */
                printf( "Pecas produzidas: " );
                scanf( "%d", &pecas );
                printf( "Valor por peca: " );
                scanf( "%f", &valorPeca );
                pagamento = pecas * valorPeca;
                break;

            default:
                printf( "Codigo invalido!\n" );
                pagamento = 0.0f;
                break;
        }

        printf( "Pagamento semanal: R$ %.2f\n\n", pagamento );

        printf( "Proximo codigo (-1 para terminar): " );
        scanf( "%d", &codigo );
    }

    return 0;
} /* fim da funcao main */
\end{lstlisting}
\end{respostabox}

\newpage

% ============================================================
\section{Exercício 4.29 -- Leis de De Morgan}
% ============================================================

\begin{respostabox}
\textbf{As Leis de De Morgan:}
\begin{align*}
  !(\text{C1} \,\&\&\, \text{C2}) &\equiv (!\text{C1} \,||\, !\text{C2}) \\
  !(\text{C1} \,||\, \text{C2}) &\equiv (!\text{C1} \,\&\&\, !\text{C2})
\end{align*}

\begin{enumerate}[label=\textbf{\alph*)}]

\item Original: \texttt{!(x < 5) \&\& !(y >= 7)}\\
Equivalente: \texttt{!(x < 5 || y >= 7)} \quad ou \quad \texttt{x >= 5 \&\& y < 7}

\item Original: \texttt{!(a == b) || !(g != 5)}\\
Equivalente: \texttt{!(a == b \&\& g != 5)} \quad ou \quad \texttt{a != b || g == 5}

\item Original: \texttt{!((x <= 8) \&\& (y > 4))}\\
Equivalente: \texttt{!(x <= 8) || !(y > 4)} \quad ou \quad \texttt{x > 8 || y <= 4}

\item Original: \texttt{!((i > 4) || (j <= 6))}\\
Equivalente: \texttt{!(i > 4) \&\& !(j <= 6)} \quad ou \quad \texttt{i <= 4 \&\& j > 6}

\end{enumerate}

\textbf{Programa de verificação:}
\begin{lstlisting}
/* Ex4.29: de_morgan.c
   Demonstra as Leis de De Morgan */
#include <stdio.h>

int main( void )
{
    int x, y;

    printf( "Digite x e y: " );
    scanf( "%d%d", &x, &y );

    printf( "\nItem (a):\n" );
    printf( "  !(x<5) && !(y>=7) = %d\n", !(x<5) && !(y>=7) );
    printf( "  x>=5 && y<7       = %d\n", x>=5 && y<7 );

    /* Repetir para os demais itens */

    return 0;
} /* fim da funcao main */
\end{lstlisting}
\end{respostabox}

% ============================================================
\section{Exercício 4.31 -- Losango}
% ============================================================

\begin{respostabox}
\begin{lstlisting}
/* Ex4.31: losango.c
   Imprime forma de losango com altura fixa */
#include <stdio.h>

int main( void )
{
    int linha, j;

    /* Metade superior: 1, 3, 5, 7, 9 asteriscos */
    for ( linha = 1; linha <= 5; ++linha ) {
        for ( j = 1; j <= 5 - linha; ++j ) {
            printf( " " );
        }
        for ( j = 1; j <= 2 * linha - 1; ++j ) {
            printf( "*" );
        }
        printf( "\n" );
    }

    /* Metade inferior: 7, 5, 3, 1 asteriscos */
    for ( linha = 4; linha >= 1; --linha ) {
        for ( j = 1; j <= 5 - linha; ++j ) {
            printf( " " );
        }
        for ( j = 1; j <= 2 * linha - 1; ++j ) {
            printf( "*" );
        }
        printf( "\n" );
    }

    return 0;
} /* fim da funcao main */
\end{lstlisting}
\end{respostabox}

\newpage

% ============================================================
\section{Exercício 4.32 -- Losango com Tamanho Variável}
% ============================================================

\begin{respostabox}
\begin{lstlisting}
/* Ex4.32: losango_variavel.c
   Imprime losango cujo numero de linhas e fornecido pelo usuario */
#include <stdio.h>

int main( void )
{
    int n;          /* numero impar (1 a 19): linhas no meio */
    int meio;       /* (n+1)/2: numero da linha do meio      */
    int linha, j;

    do {
        printf( "Digite um numero impar (1 a 19): " );
        scanf( "%d", &n );
    } while ( n < 1 || n > 19 || n % 2 == 0 );

    meio = ( n + 1 ) / 2;

    /* Metade superior */
    for ( linha = 1; linha <= meio; ++linha ) {
        for ( j = 1; j <= meio - linha; ++j ) {
            printf( " " );
        }
        for ( j = 1; j <= 2 * linha - 1; ++j ) {
            printf( "*" );
        }
        printf( "\n" );
    }

    /* Metade inferior */
    for ( linha = meio - 1; linha >= 1; --linha ) {
        for ( j = 1; j <= meio - linha; ++j ) {
            printf( " " );
        }
        for ( j = 1; j <= 2 * linha - 1; ++j ) {
            printf( "*" );
        }
        printf( "\n" );
    }

    return 0;
} /* fim da funcao main */
\end{lstlisting}
\end{respostabox}

% ============================================================
\section{Exercício 4.33 -- Equivalentes Romanos de 1 a 100}
% ============================================================

\begin{respostabox}
\textbf{Algoritmo:} usar repetidamente a tabela de valores romanos: I=1, IV=4, V=5,
IX=9, X=10, XL=40, L=50, XC=90, C=100. Subtraímos o maior valor possível enquanto
adicionamos o símbolo correspondente.

\begin{lstlisting}
/* Ex4.33: romanos.c
   Tabela de equivalentes em algarismos romanos de 1 a 100 */
#include <stdio.h>

int main( void )
{
    int n, valor;

    printf( "%5s   %s\n", "Dec", "Romano" );

    for ( n = 1; n <= 100; ++n ) {
        printf( "%5d   ", n );
        valor = n;

        /* C = 100 */
        while ( valor >= 100 ) { printf( "C" ); valor -= 100; }
        /* XC = 90 */
        while ( valor >= 90 )  { printf( "XC" ); valor -= 90; }
        /* L = 50 */
        while ( valor >= 50 )  { printf( "L" ); valor -= 50; }
        /* XL = 40 */
        while ( valor >= 40 )  { printf( "XL" ); valor -= 40; }
        /* X = 10 */
        while ( valor >= 10 )  { printf( "X" ); valor -= 10; }
        /* IX = 9 */
        while ( valor >= 9 )   { printf( "IX" ); valor -= 9; }
        /* V = 5 */
        while ( valor >= 5 )   { printf( "V" ); valor -= 5; }
        /* IV = 4 */
        while ( valor >= 4 )   { printf( "IV" ); valor -= 4; }
        /* I = 1 */
        while ( valor >= 1 )   { printf( "I" ); valor -= 1; }

        printf( "\n" );
    }

    return 0;
} /* fim da funcao main */
\end{lstlisting}

\textbf{Exemplos:} 14 = XIV, 49 = XLIX, 99 = XCIX, 100 = C.
\end{respostabox}

\newpage

% ============================================================
\section{Exercício 4.34 -- \texttt{do...while} vs.\ \texttt{while}}
% ============================================================

\begin{respostabox}
\textbf{Substituir \texttt{do...while} por \texttt{while}:}

A diferença essencial entre as duas estruturas é que \texttt{do...while} \textbf{sempre
executa o corpo pelo menos uma vez}, enquanto \texttt{while} pode não executar nenhuma
vez (se a condição já for falsa de início).

\textbf{Para converter \texttt{do...while} \(\to\) \texttt{while}:} basta executar o
corpo \textit{antes} do laço:
\begin{lstlisting}
/* Original: do-while */
do {
    /* corpo */
} while ( condicao );

/* Equivalente com while */
/* corpo */     /* primeira execucao garantida */
while ( condicao ) {
    /* corpo */ /* repeticao se condicao verdadeira */
}
\end{lstlisting}

\textbf{Problema ao converter \texttt{while} \(\to\) \texttt{do...while}:} se a
condição já for falsa no início, o \texttt{while} \textbf{não executa o corpo nem uma
vez}, mas o \texttt{do...while} \textbf{executa pelo menos uma vez}, o que pode causar
comportamento incorreto.

\textbf{Solução:} use uma instrução \texttt{if} para envolver o \texttt{do...while}:
\begin{lstlisting}
if ( condicao ) {
    do {
        /* corpo */
    } while ( condicao );
}
\end{lstlisting}

Assim, se a condição já for falsa, nem o \texttt{if} nem o \texttt{do...while} executam.
\end{respostabox}

% ============================================================
\section{Exercício 4.35 -- Removendo \texttt{break}}
% ============================================================

\begin{respostabox}
\textbf{Como remover \texttt{break} de um laço:}

A ideia central é introduzir uma \textbf{variável de controle adicional} (uma ``flag'')
que indique se o laço deve continuar. A condição original é estendida com essa flag,
unindo as duas formas de saída em uma só (na condição de continuação).

\textbf{Padrão original com \texttt{break}:}
\begin{lstlisting}
while ( condicao_principal ) {
    /* ... */
    if ( condicao_de_saida ) {
        break;
    }
    /* ... */
}
\end{lstlisting}

\textbf{Equivalente sem \texttt{break}:}
\begin{lstlisting}
int continuar = 1;  /* flag de continuacao */
while ( condicao_principal && continuar ) {
    /* ... */
    if ( condicao_de_saida ) {
        continuar = 0;  /* sinaliza saida; loop terminara apos */
    }
    else {
        /* ... codigo apos o break original ... */
    }
}
\end{lstlisting}

A divisão do código em \texttt{if/else} garante que, ao detectar a saída, o restante
do corpo \textit{não} seja executado nessa iteração. A condição estendida garante que
o laço termina na próxima checagem.
\end{respostabox}

\newpage

% ============================================================
\section{Exercício 4.36 -- O que faz o segmento de programa?}
% ============================================================

\begin{respostabox}
\begin{lstlisting}
for ( i = 1; i <= 5; i++ ) {
    for ( j = 1; j <= 3; j++ ) {
        for ( k = 1; k <= 4; k++ )
            printf( "*" );
        printf( "\n" );
    }
    printf( "\n" );
}
\end{lstlisting}

\textbf{Análise dos três loops aninhados:}
\begin{itemize}[noitemsep]
  \item Loop \texttt{i} (mais externo): executa 5 vezes
  \item Loop \texttt{j} (médio): para cada \texttt{i}, executa 3 vezes
  \item Loop \texttt{k} (mais interno): para cada \texttt{j}, imprime 4 asteriscos
\end{itemize}

Para cada iteração de \texttt{i}, são impressas 3 linhas com 4 asteriscos cada
(forma um retângulo $3 \times 4$), seguidas de uma linha em branco.

\textbf{Saída:} \textbf{5 grupos} de retângulos $3 \times 4$ de asteriscos, separados
por linhas em branco. Total: $5 \times 3 = 15$ linhas com 4 asteriscos, mais 5 linhas
em branco.

\begin{saida}
****\\****\\****\\
\\
****\\****\\****\\
\\
****\\****\\****\\
\\
****\\****\\****\\
\\
****\\****\\****
\end{saida}
\end{respostabox}

% ============================================================
\section{Exercício 4.37 -- Removendo \texttt{continue}}
% ============================================================

\begin{respostabox}
\textbf{Como remover \texttt{continue} de um laço:}

O comando \texttt{continue} pula o restante da iteração atual e prossegue para a
próxima. Para removê-lo, basta colocar o restante do código do laço dentro de um
\texttt{else} (ou um \texttt{if} com a condição negada).

\textbf{Padrão original com \texttt{continue}:}
\begin{lstlisting}
while ( condicao ) {
    /* parte 1 */
    if ( condicao_pular ) {
        continue;
    }
    /* parte 2 */
}
\end{lstlisting}

\textbf{Equivalente sem \texttt{continue}:}
\begin{lstlisting}
while ( condicao ) {
    /* parte 1 */
    if ( !condicao_pular ) {  /* condicao negada */
        /* parte 2 */
    }
}
\end{lstlisting}

\textbf{Aplicação à Figura 4.12 (programa de soma de 1 a 10 pulando 5):}
\begin{lstlisting}
/* Original com continue: */
for ( x = 1; x <= 10; ++x ) {
    if ( x == 5 ) {
        continue;
    }
    printf( "%d ", x );
}

/* Equivalente sem continue: */
for ( x = 1; x <= 10; ++x ) {
    if ( x != 5 ) {  /* condicao negada */
        printf( "%d ", x );
    }
}
\end{lstlisting}
\end{respostabox}

\end{document}
